% ----- CHAUPAI 9 -----

\newpage

\subsection*{\deva{नवम चौपाई} (Ninth Chaupai)}
\addcontentsline{toc}{subsection}{\deva{नवम चौपाई} (Ninth Chaupai) — \deva{सूक्ष्म रूप...}}

\begin{minipage}[t]{0.55\textwidth}
\vspace{0pt}

{\large \linespread{1.5}\selectfont
\deva{सूक्ष्म रूप धरि सियहिं दिखावा।}\\
\deva{बिकट रूप धरि लंक जरावा॥}
\par}

\vspace{0.5em}

{\large \linespread{1.5}\selectfont
\telu{సూక్ష్మ రూప ధరి సియహిఁ దిఖావా।}\\
\telu{బికట రూప ధరి లంక జరావా॥}
\par}

\vspace{0.5em}

{\large \linespread{1.3}\selectfont
\textit{sūkṣma rūpa dhari siyahiṃ dikhāvā |}\\
\textit{bikaṭa rūpa dhari laṃka jarāvā ||}
\par}
\end{minipage}%
\hfill
\begin{minipage}[t]{0.42\textwidth}
\vspace{0pt}
\begin{tcolorbox}[anchorbox]
\textbf{Emotional Intelligence and Scale:} Hanuman knew exactly how to scale his approach. He approached allies in a tiny, unthreatening form (\textit{Sukshma Rupa}) to build trust and offer comfort, but assumed a massive, formidable form (\textit{Bikata Rupa}) when he needed to defend himself. Emotional intelligence requires adapting your approach to match the needs of the situation.
\vspace{0.5em}\newline He made himself intentionally unthreatening to sneak into Lanka, showing that true strength knows when to be subtle.\newline\null\hfill\href{https://www.valmiki.iitk.ac.in/sloka?field_kanda_tid=5&language=dv&field_sarga_value=2}{\textit{Sundara Kanda, 2:49}}
\end{tcolorbox}
\end{minipage}

\vspace{0.5em}
\subsubsection*{\textbf{Visual Rhythm Map:}}
\begin{table}[h!]
\centering
\renewcommand{\arraystretch}{1.2}
\setlength{\tabcolsep}{0pt}
\begin{tabularx}{\textwidth}{|*{16}{>{\centering\arraybackslash\hsize=1.0\hsize}X|}}
\hline
\multicolumn{3}{|>{\centering\arraybackslash\hsize=3.0\hsize}X|}{\textbf{\guru{सू}\laghu{क्ष्म}} \par (3)} &
\multicolumn{3}{>{\centering\arraybackslash\hsize=3.0\hsize}X|}{\textbf{\guru{रू}\laghu{प}} \par (3)} &
\multicolumn{2}{>{\centering\arraybackslash\hsize=2.0\hsize}X|}{\textbf{\laghu{ध}\laghu{रि}} \par (2)} &
\multicolumn{3}{>{\centering\arraybackslash\hsize=3.0\hsize}X|}{\textbf{\laghu{सि}\laghu{य}\laghu{हिँ}} \par (3)} &
\multicolumn{5}{>{\centering\arraybackslash\hsize=5.0\hsize}X|}{\textbf{\laghu{दि}\guru{खा}\guru{वा}} \par (5)} \\
\hline
\end{tabularx}
\vspace{-1.5em}
\end{table}

\begin{table}[h!]
\centering
\renewcommand{\arraystretch}{1.2}
\setlength{\tabcolsep}{0pt}
\begin{tabularx}{\textwidth}{|*{16}{>{\centering\arraybackslash\hsize=1.0\hsize}X|}}
\hline
\multicolumn{3}{|>{\centering\arraybackslash\hsize=3.0\hsize}X|}{\textbf{\laghu{बि}\laghu{क}\laghu{ट}} \par (3)} &
\multicolumn{3}{>{\centering\arraybackslash\hsize=3.0\hsize}X|}{\textbf{\guru{रू}\laghu{प}} \par (3)} &
\multicolumn{2}{>{\centering\arraybackslash\hsize=2.0\hsize}X|}{\textbf{\laghu{ध}\laghu{रि}} \par (2)} &
\multicolumn{3}{>{\centering\arraybackslash\hsize=3.0\hsize}X|}{\textbf{\guru{लं}\laghu{क}} \par (3)} &
\multicolumn{5}{>{\centering\arraybackslash\hsize=5.0\hsize}X|}{\textbf{\laghu{ज}\guru{रा}\guru{वा}} \par (5)} \\
\hline
\end{tabularx}
\vspace{-1.5em}
\end{table}

\vspace{1em}
\textbf{ \deva{सरल-संस्कृतम्}:}\\
 \deva{भवान् सूक्ष्मं (लघु) रूपं धृत्वा सीतायै (सियहिं) दर्शनं दत्तवान्।}\\
 \deva{भवान् विकटं (भयानकं/विशालं) रूपं धृत्वा लङ्कां भस्मीकृतवान् (जरावा)।}\\


Assuming a microscopic form, you revealed yourself to Mother Sita.\\
Assuming a massive/terrifying form, you burned down Lanka.


\vspace{1em}
\newpage
\textbf{\deva{प्रतिपदार्थः} (Meaning of each word):}
\begin{table}[h!]
\centering
\renewcommand{\arraystretch}{1.2}
\begin{tabularx}{\textwidth}{@{} l l X X @{}}
\toprule
\textbf{Awadhi} & \textbf{Sanskrit} & \textbf{English} & \textbf{Grammar \& Notes} \\
\midrule
\deva{सूक्ष्म रूप} / \telu{సూక్ష్మ రూప}  &  \deva{सूक्ष्म-रूपम्}  &  Microscopic form  &   \\
\deva{धरि} / \telu{ధరి}  &  \deva{धृत्वा}  &  Having assumed  &   \\
\deva{सियहिं} \gotchaicon / \telu{సియహిఁ}  &  \deva{सीतायै}  &  To Mother Sita  & Dative case suffix -hiṃ \\
\deva{दिखावा} / \telu{దిఖావా}  &  \deva{दर्शितवान्}  &  Showed / Revealed  &   \\
\deva{बिकट रूप} / \telu{బికట రూప}  &  \deva{विकट-रूपम्}  &  Terrifying form  & \\ 
\deva{धरि} / \telu{ధరి}  &  \deva{धृत्वा}  &  Having assumed  &   \\
\deva{लंक} / \telu{లంక}  &  \deva{लङ्काम्}  &  Lanka  &   \\
\deva{जरावा} / \telu{జరావా}  &  \deva{ज्वालितवान्}  &  Burned  & Jval $\rightarrow$ Jaraavaa \\ 
\bottomrule
\end{tabularx}
\end{table}

\begin{tcolorbox}[gotchabox]
\begin{itemize}
    \item \textbf{The Dative Suffix (\deva{सियहिं}):} In Awadhi, the formal name \deva{सीता} softens to \deva{सिय} (Siya). The \deva{-हिं} (-hiṃ) is the objective marker we saw earlier in \deva{मोहिं} (to me). So \deva{सियहिं} means "To Mother Siya".
    \item \textbf{Phonetic Shifts:} Sanskrit \deva{विकट} (Vikata) becomes \deva{बिकट} (Bikata) via the $v \rightarrow b$ shift. Also, the causative verb \deva{जरावा} (Jaraavaa) comes from the root \deva{ज्वल्} (to burn), meaning "caused to burn".
\end{itemize}
\end{tcolorbox}



\newpage
